\chapter{Многомасштабность, слои и забывание}
\label{ch:timescales}

\section{Протекающий интегратор: точный смысл}

Память Рин расслоена по временным масштабам: \emph{сессия} $\to$ \emph{день} $\to$
\emph{сезон} $\to$ \emph{год} $\to$ \emph{жизнь}. Каждый слой — \emph{протекающий
интегратор} (leaky integrator): состояние, которое накапливает вход и понемногу
«протекает» (теряет часть себя) на каждом шаге.

\begin{definition}[Протекающий интегратор]
Состояние $s$, обновляемое под входом $x_t$ по правилу
\begin{equation}
s_t \;=\; \underbrace{\Big(1-\tfrac{1}{\tau}\Big)}_{\text{удержание}}\,s_{t-1}
\;+\; \underbrace{\tfrac{1}{\tau}}_{\text{впитывание}}\,x_t,
\qquad \tau>1.
\label{eq:leaky}
\end{equation}
Коэффициент $(1-1/\tau)$ — доля прошлого состояния, которая \emph{удерживается};
$1/\tau$ — доля нового входа, которая \emph{впитывается}. Постоянная времени $\tau$
задаёт «медленность»: большое $\tau$ — слой почти не меняется за шаг (удерживает
долго, впитывает по чуть-чуть), малое $\tau$ — резко реагирует и быстро забывает.
\end{definition}

\begin{intuition}[Ведро с дыркой]
Протекающий интегратор — это ведро с дыркой в дне, куда капает вода (вход). Уровень
воды $s$ растёт от капель и падает от протечки. Широкая дырка (малое $\tau$) —
уровень скачет за каждой каплей и быстро сходит; узкая дырка (большое $\tau$) —
уровень меняется медленно, зато держится долго. «Сессия» — ведро с широкой дыркой,
«жизнь» — с почти закрытой.
\end{intuition}

\section{Два разных смысла «слоёв» — не путать}
\label{sec:twolayers}

Здесь важная развилка, которую легко замылить: слово «слои» и «протекание» означают
\emph{разное} для настроения и для памяти.

\paragraph{Настроение — параллельные интеграторы, переноса нет.}
Слои настроения (глава~\ref{ch:mood}) — это несколько интеграторов
\eqref{eq:leaky} с разными $\tau_\ell$, работающих \emph{параллельно на одном и том
же потоке событий}. Старший слой не питается от младшего; оба питаются событиями,
просто старший интегрирует медленнее. «Младший дрожит сильнее» — это высокочастотная
компонента того же сигнала. \textbf{Переноса содержания между слоями нет} — есть
разночастотная фильтрация одного входа, а актуальное настроение — суперпозиция.

\paragraph{Память — перенос есть, это повышение слоя консолидацией.}
Слои памяти устроены иначе: узел \emph{повышается} из быстрого слоя в медленный, когда
консолидация делает его устойчивым (раздел~\ref{sec:consolidation}). Вот это и есть
настоящее «протекание младшего в старший»: не диффузия сигнала, а
\emph{перемаркировка} узла (\texttt{layer}: session $\to$ day $\to \dots$) при
накоплении силы. Механизм известен по форме; конкретные пороги повышения — открыты.

\begin{openq}[Пороги повышения слоя]
При какой накопленной силе / частоте обращений / значимости узел повышается в
следующий слой? Разумные кандидаты: порог по базовой активации ACT-R
(раздел~\ref{sec:decay}); порог по числу независимых подтверждений; порог по
$|\conf|$. Точные значения и то, единый ли это порог для всех типов узлов, — пока не
зафиксированы; это параметр для стенда (глава~\ref{ch:experiment}).
\end{openq}

\section{Консолидация как огрубление графа}
\label{sec:consolidation}

Переход от эпизодов к смыслам — это \emph{огрубление} (coarsening): десять записей
«зритель заходил в понедельник», «\ldots в среду», \ldots заменяются одним узлом
«зритель — частый гость». Мелкая структура схлопывается в крупную; часть информации
теряется намеренно.

Как понять, что огрубление \emph{хорошее}? Здесь мост к математике DTI. Введём:
\begin{itemize}\itemsep1pt
\item $L_f$ — лапласиан \emph{мелкого} (fine) графа; лапласиан описывает, как сигнал
диффундирует по графу (кто на кого влияет и насколько);
\item $L_c$ — лапласиан \emph{крупного} (coarse) графа после огрубления;
\item $P$ — оператор \emph{проекции}: сопоставляет каждому мелкому узлу его крупную
группу (например, десять эпизодов $\mapsto$ узел «частый гость»);
\item $\Phi(X,L)$ — оператор распространения (одна или несколько итераций передачи
сообщений) сигнала $X$ по графу с лапласианом $L$.
\end{itemize}
Хорошее огрубление \emph{коммутирует} с распространением:
\begin{equation}
\big\|\,\underbrace{P\,\Phi_f(X,L_f)}_{\text{распространить, потом огрубить}}
\;-\;\underbrace{\Phi_c(PX,L_c)}_{\text{огрубить, потом распространить}}\,\big\|
\;\le\;\varepsilon .
\label{eq:commute}
\end{equation}
Читается так: неважно, в каком порядке действовать — сначала «подумать» на мелком
графе и затем сжать, или сначала сжать и «подумать» на крупном, — результат должен
почти совпасть. Норма $\|\cdot\|$ берёт разницу по всем входным сигналам $X$;
$\varepsilon$ — мера потери. Малое $\varepsilon$ = огрубление сохранило то, что нужно
для рассуждения; большое = сжатие сломало динамику.

\begin{workedbox}[Что должно выжить при консолидации]
Свернём десять эпизодов визитов в узел «частый гость». \emph{Выжить} обязаны:
частота (была высокой), устойчивая валентность (тёплая), тип связи (регулярность).
\emph{Потеряться} можно: точные даты, дословные реплики. Плохое огрубление сохранит
даты и потеряет валентность — тогда Рин будет помнить \emph{когда} заходил зритель,
но забудет, \emph{как относится} к нему; это большая $\varepsilon$ по аффективной
оси. Значит критерий огрубления — сохранять аффективно- и реляционно-релевантное
подпространство, а не спектр целиком.
\end{workedbox}

\section{Забывание как тень консолидации}
\label{sec:decay}

Забывание — не отдельный механизм, а \emph{отрицательная сторона} того же процесса:
усиливай используемое, ослабляй неиспользуемое.

\subsection{Правильный decay — это ACT-R}

\emph{ACT-R} (Adaptive Control of Thought—Rational) — влиятельная когнитивная
архитектура Джона Андерсона, модель того, как устроены человеческие память и
научение. Нам из неё нужен один компонент — \emph{базовая активация} (base-level
activation): сила следа в памяти есть функция \emph{частоты} и \emph{свежести}
обращений со степенным затуханием:
\begin{equation}
B_v \;=\; \ln\!\Big(\sum_{k=1}^{n} t_k^{-\alpha}\Big).
\label{eq:actr}
\end{equation}
Разбор: $t_k$ — сколько времени назад было $k$-е обращение к узлу; каждое обращение
вносит вклад $t_k^{-\alpha}$, и чем оно \emph{свежее} (меньше $t_k$), тем вклад
больше; сумма по всем $n$ обращениям даёт накопленную силу (\emph{частота});
логарифм сжимает шкалу; $\alpha>0$ — скорость забывания. Два следствия, которых нет
у плоской «$\times 0{,}95$ в день»: частые обращения укрепляют след, а
\emph{разнесённые во времени} укрепляют сильнее массовых (эффект распределённого
повторения — spacing).

\begin{intuition}[Тропинка в лесу]
Тропинку протаптывают многими проходами; если по ней ходят изредка, но регулярно,
она держится лучше, чем если пройти сто раз за день, а потом забросить. Формула
\eqref{eq:actr} именно это и кодирует: важна не только суммарная частота, но и её
распределение во времени.
\end{intuition}

\begin{pitfall}[Доступ сам тренирует память — различай access\_reason]
Тонкий баг: не всякое обращение к узлу должно укреплять его. Если ночной скраб,
debug-инспектор или retrieval-тест часто трогают узел, он не должен от этого
становиться «важным жизненным следом» — иначе \emph{сама диагностика тренирует
память}. Поэтому у каждого доступа есть \texttt{access\_reason}: \texttt{user\_recall},
\texttt{model\_recall}, \texttt{maintenance\_touch}, \texttt{debug\_read},
\texttt{scrub\_visit}, \texttt{context\_packed}. В ACT-R-активацию \eqref{eq:actr}
входят с полным весом только \emph{подлинные} припоминания (user/model recall);
обслуживающие тачи — с нулевым или почти нулевым. Наблюдатель не должен менять
наблюдаемое.
\end{pitfall}

\subsection{Прунинг: не только по использованию}

Прунинг только по использованию опасен: редкий, но критичный факт вымоется как
мусор. Поэтому склонность к удалению гейтируется \emph{важностью} узла:
\begin{equation}
\text{prune}(v)\;\propto\;\text{неиспользование}(v)\,\cdot\,\big(1-\text{importance}(v)\big).
\end{equation}
Здесь \texttt{importance} — атрибут узла (см.\ схему в главе~\ref{ch:substrate}),
задаваемый значимостью при кодировании (salience из appraisal,
глава~\ref{ch:personality}) и подкрепляемый структурной центральностью
(глава~\ref{ch:topology}). Редкое-но-важное выживает; частое-но-пустое держит
частота; вымывается лишь редкое-и-неважное. Врождённым (глубоко консолидированным)
узлам задаётся ниже скорость распада — континуум глубины, а не бинарность.

\subsection{importance нельзя диффундировать по путям}
\label{sec:importance-leak}

Здесь важная ловушка. Соблазн: раз узел-домен «цепи Маркова» важен, пусть раздаёт
$\text{importance}\to 1$ своим детям, те — своим, и так область защищена. \textbf{Так
делать нельзя.} importance — это диффузия по графу (глава~\ref{ch:topology}), а
диффузия растекается \emph{везде}: через несколько хопов по каким-нибудь рёбрам она
доползёт до «помыть машину» (которой у тебя нет), и та станет «важной». Графовые пути
неисповедимы; неуправляемая передача importance приводит к тому, что важным
становится \emph{всё}, а значит — ничто.

Правильно — \textbf{не пускать importance по произвольным путям, а закреплять её
членством в домене}:

\begin{itemize}\itemsep2pt
\item \textbf{importance по членству, а не по диффузии.} У узла есть \emph{теги
доменов}, у доменов — importance. Эффективная важность узла $=$ функция от
importance его доменов $+$ собственная salience. Членство назначается \emph{при
кодировании} (прочитали учебник $\to$ его узлы помечены \texttt{domain=Markov},
высокая importance), а не распространяется обходом. «Помыть машину» никогда не
получает тег \texttt{Markov}, поэтому \emph{ничего не наследует}. Утечка исключена по
построению.
\item \textbf{importance по намерению.} То, что \emph{попросили} выучить, получает
высокую importance по провенансу приобретения (в отличие от случайного тривиа).
\item \textbf{Если всё же нужна передача} (для непомеченных узлов) — только
\emph{типо-гейтированная} (по рёбрам членства \textsc{isA}/\textsc{partOf}, но не по
ассоциативным), \emph{сильно затухающая} ($\lambda^k$ с малым $\lambda$) и
\emph{ограниченная} 1--2 хопами; и это надо \emph{проверять} на стенде, что importance
осталась локальной.
\end{itemize}

Тогда защита области (глава~\ref{ch:runtime}, раздел~\ref{sec:storage}) работает
аккуратно: домен-хаб с высокой importance и низким распадом $+$ члены, наследующие
\emph{по членству}; детали могут тускнеть, но gist области держится, а сон-репетиция
периодически освежает высоко-importance домены — без утечки на «помыть машину».

\section{Что процесс делает с найденной конфигурацией}
\label{sec:whatfusiondoes}

Предположим, фоновый разведчик (глава~\ref{sec:background}) со-активировал тройку
$\{A,B,C\}$, которая выглядит связной. Что происходит дальше? Разделю на
\emph{известное} и \emph{открытое}, чтобы не размывать.

\paragraph{Что известно (устоявшийся конвейер).}
\begin{enumerate}\itemsep1pt
\item \textbf{Композиция.} Языковая модель по объединению узлов и их контекстных
ореолов (глава~\ref{ch:fusion}) формулирует \emph{кандидата} $C^\ast$ — обобщение
или вывод.
\item \textbf{Провизорная запись.} $C^\ast$ пишется с низкой уверенностью и
провенансом $\{A,B,C\}$; на различающих пробелах ставится \texttt{need\_to\_discover}.
\item \textbf{Конкуренция.} Дальше $C^\ast$ живёт по общим правилам: используется и
корроборируется $\Rightarrow$ сила растёт, узел повышается в слой
(раздел~\ref{sec:twolayers}); игнорируется $\Rightarrow$ распадается
(раздел~\ref{sec:decay}); опровергается $\Rightarrow$ понижается по TMS
(глава~\ref{ch:epistemics}).
\end{enumerate}

\paragraph{Что открыто.} Как \emph{решить}, что со-активация $\{A,B,C\}$ достаточно
связна, чтобы вообще пытаться её сшить (а не является случайной кучей), — не решено.

\begin{openq}[Распознавание связности со-активации]
Кандидаты: \emph{(а)} LLM-судья («есть ли тут осмысленный вывод?») — гибко, но не
доказуемо и стоит вызова модели; \emph{(б)} графовые шаблоны-схемы («потребность $+$
ограничение $+$ окно $\to$ перенести действие») — дёшево, но покрывают лишь частые
формы; \emph{(в)} обучаемый скорер связности — требует разметки и вносит
непрозрачность. Ни один пока не признан удачным; на старте — (а) с фильтром $\varphi$
и обратной связью (главы~\ref{ch:personality},~\ref{ch:voice}). Это и есть «трудные
80\%» из главы~\ref{ch:fusion}.
\end{openq}

\subsection{Забывать ли посылки после fusion?}

Возникает соблазн: раз $\{A,B\}\to C$ состоялось, посылки $A,B$ больше не нужны —
удалить. Это \textbf{неверно} (слишком мануально и опасно). Правильно — не удалять
насильно, а положиться на естественный распад:

\begin{itemize}\itemsep1pt
\item при \emph{хорошей} консолидации к $A,B$ просто перестают обращаться — вся
активность идёт к более полезному $C$; базовая активация \eqref{eq:actr} $A,B$
падает, и они отмирают \emph{сами} (важные — выживают). Забывание тут
\emph{эмерджентно}, а не по правилу «удалить родителей»;
\item удалять $A,B$ немедленно нельзя ещё и потому, что они — \emph{провенанс} $C$:
если $C$ позже опровергнут, по $A,B$ его пересматривают или переводят
(глава~\ref{ch:epistemics}). Убрав посылки, теряешь возможность ревизии.
\end{itemize}
Итог по-человечески верный: детали тускнеют, суть остаётся, но след обоснования
хранится, пока не осыпался по неиспользованию.

\begin{principle}[support skeleton: обоснование переживает распад деталей]
Здесь есть скрытый конфликт: провенанс требует держать посылки для будущей ревизии
(TMS), но детали естественно \emph{decay}-ятся. Разрешение — не хранить полный
богатый эпизод вечно, а при распаде оставлять \emph{support skeleton}: минимальный
стаб \{id, source, time, hash, node\_kind, claim-summary\}, достаточный для ретракции.
Тогда либо-либо снимается: подробность уходит, а \emph{способность ревизовать}
остаётся. Без скелета либо память не забывает, либо потом нечем пересматривать
выводы.
\end{principle}
